%% notes-tex/031-yeast-chemical-space/sections/3-published.tex

\section{The expansion has already been done\secstatus{todo}}
\label{sec:published}

The scoping question has a published answer, and it arrived in 2026. Yeast-MetaTwin
\external{Wu et al.\ 2026, ``Systematically exploring yeast metabolism through
retrobiosynthesis and deep learning''; Nielsen and Kerkhoven are authors} applies exactly the
method under consideration here, at a scale that would be hard to improve on casually.

The numbers, quoted through \file{yeast9\_literature\_counts.py} and plotted hatched in
Figure~\ref{fig:scope}. The Yeast Metabolome Database holds \textbf{16{,}042} entries carrying
a SMILES, of which \textbf{14{,}882 are absent from Yeast9}, and those missing entries are
\textbf{14{,}310 lipids against 572 non-lipids}. The authors extracted \textbf{21{,}921
enzyme-associated reaction rules and 213 spontaneous rules} from MetaNetX and MetaCyc, report
that 95 percent of the rules regenerate their original reactions, and trimmed the resulting
pool to a connected network of \textbf{1{,}092{,}946 reactions}. The outcome they state is
that coverage of the yeast metabolome rose \textbf{from 7 percent to 92 percent}, and for the
non-lipid metabolome from 54 to 75 percent.

%% SOURCE: experiments/031-env-chemgen-inhibitor-tolerance/scripts/yeast9_chemical_accounting.py
%% and yeast9_literature_counts.py -> yeast9_scope.svg, converted by `make plots`.
\tcfig{figures/yeast9_scope.pdf}{\textbf{The retrobiosynthetic
expansion is published, and what is left open is stereochemistry, lipids and a missing mass
fraction.} Hatched bars are published and were not measured here; solid bars are measured from
the release. \textbf{a}, Metabolome coverage before and after the published expansion, for the
whole metabolome and for the non-lipid subset. \textbf{b}, Yeast dry-mass composition computed from the
model's own biomass pseudo-reactions, which closes to 96.1 percent; red whiskers are the
published range for the four pools the literature covers. DNA is 0.39 percent and the cofactor
and ion pools together 0.73 percent, so the remainder after the four large pools is not a
five-to-ten percent small-molecule fraction. \textbf{c}, What remains: Yeast9 species whose structure names a set
rather than a molecule, Yeast9 lipid species with no structure at all, and the non-lipid
metabolites the published expansion could not connect.}{fig:scope}

Two consequences follow immediately. First, proposing a retrobiosynthetic enumeration of yeast
metabolism as new work would duplicate a published result from the same group that maintains
the reconstruction. Second, the residual they report is specific and small: \textbf{267
non-lipid metabolites}, 1.7 percent of the metabolome, could not be connected without
introducing metabolites outside the yeast metabolome.

\notesec{The mass question, answered from the model's own biomass equation}
The rounded figures in the literature do not add up and say nothing about DNA or the
small-molecule pool: yeast is reported as 40 to 50 percent protein, about 10 percent RNA,
about 10 percent lipid and about 25 percent cell wall \external{Milo and Phillips 2016;
Feldmann 2012}. The reconstruction carries a complete answer that needs no lookup. Its biomass
pseudo-reaction draws one unit each from a protein, carbohydrate, RNA, DNA, lipid, cofactor and
ion pool, and every pool lists its constituents with coefficients in millimoles per gram dry
weight, so multiplying by formula mass gives grams per gram directly
(Figure~\ref{fig:scope}b).

Measured that way the pools are \textbf{46.8 percent protein, 38.3 percent carbohydrate, 6.4
percent RNA and 3.5 percent lipid}, with \textbf{0.48 percent cofactor, 0.39 percent DNA and
0.25 percent ion}. They sum to 96.1 percent, and the 3.9 percent shortfall is the few
constituents carrying no formula, so the accounting nearly closes and the residual is named
rather than left implicit.

That settles the question the rounded figures leave open, and the answer is not what the
arithmetic suggests. \textbf{The remainder after protein, carbohydrate, RNA and lipid is not
five to ten percent of small molecules.} DNA is 0.39 percent and the cofactor and ion pools
together are 0.73 percent. Two caveats on reading those numbers. The model's carbohydrate pool
runs above the 25 percent cell-wall figure because it also carries storage glycogen and
trehalose, and its lipid pool runs below the published 10 percent. And the cofactor and ion
pools are a \textbf{lower bound} on the small-molecule pool, because they hold the cofactors
and ions the biomass equation requires rather than the cell's whole metabolite pool; a soluble
pool measured by extraction would be larger, and the \org{E. coli} figure for that quantity is
3 to 3.9 percent of dry weight.

The practical consequence for a chemical accounting is that the pool it is about is small. If
the target is every distinct chemical species, the species are numerous and they are under one
percent of the cell by mass, so the exercise is about identity and connectivity rather than
about mass balance.
